\documentclass[../main.tex]{subfiles}
\graphicspath{{\subfix{../images/funding_rate/}}}

\begin{document}

\subsection{Funding Rate Arbitrage: Arbitraje de tasas de financiamiento}
Esta estrategia se basa en capturar el \emph{funding rate} utilizando contratos perpetuos de Binance o Bybit (por ejemplo, BTCUSDT Perpetual) en combinación con Quantfury.

\subsubsection{¿Qué es el Funding Rate?}
En los futuros perpetuos no existe una fecha de vencimiento. Para mantener el precio del contrato cercano al del mercado spot, los exchanges aplican una tasa de financiamiento, conocida como \emph{funding rate}, que se paga entre traders cada 8 horas, según el desequilibrio entre posiciones largas y cortas:

\begin{itemize}
    \item Funding rate positivo: los largos pagan a los cortos.
    \item Funding rate negativo: los cortos pagan a los largos.
\end{itemize}

\subsubsection{Objetivo}
El objetivo de la estrategia es capturar el \emph{funding rate} sin asumir riesgo direccional del mercado. Esto se logra mediante una posición \emph{hedgeada}:

\begin{quote}
Abrir simultáneamente una posición larga en el mercado spot (o un derivado que lo replique) y una posición corta en el contrato perpetuo.
\end{quote}

\subsubsection{Estadísticas de Bybit}
La tasa de financiación típica es de $0.01\%$ cada 8 horas (valor estándar), equivalente a $0.03\%$ diario.  
Suponiendo una tasa constante de $0.01\%$ por sesión, la tasa anual aproximada sería:
\[
\text{TasaAnual} = 0.01\% \cdot 3 \cdot 365 \approx 10.95\%
\]

El promedio histórico desde el lanzamiento de BTCUSDT Futures es:
\[
\text{mean} = 0.0138\%
\]
Este valor es ligeramente superior al estándar de $0.01\%$.

\begin{figure}[htbp]
    \centering
    \includegraphics[width=0.8\linewidth]{./images/funding_rate/heatmap_month_year.png}
    \caption{Heatmap del Funding Rate por mes y año (2020-2025).}
    \label{fig:heatmap_funding_rate}
\end{figure}

\begin{figure}[htbp]
    \centering
    \includegraphics[width=0.8\linewidth]{./images/funding_rate/total_year.png}
    \caption{Porcentaje total de Funding Rate por año.}
    \label{fig:total_year}
\end{figure}

\subsubsection{Ejemplo práctico}
Supongamos un capital total de $2\,000\,\text{USDT}$, distribuido de la siguiente manera:

\begin{itemize}
    \item $1\,000\,\text{USDT}$ en Quantfury.
    \item $1\,000\,\text{USDT}$ en Bybit/Binance Futures.
\end{itemize}

En Quantfury se dispone de un apalancamiento de $20\times$, lo que permite abrir una posición equivalente a $20\,000\,\text{USDT}$. Se utiliza el $40\%$ de este poder de trading ($8\,000\,\text{USDT}$) para la estrategia.  

Los pasos a seguir son:

\begin{itemize}
    \item Abrir una posición \textbf{LONG} en Quantfury por $8\,000\,\text{USDT}$.
    \item Abrir una posición \textbf{SHORT} en Bybit/Binance Futures por $8\,000\,\text{USDT}$.
\end{itemize}

\begin{figure}[htbp]
    \centering
    \includegraphics[width=0.8\linewidth]{./images/funding_rate/example_heatmap.png}
    \caption{Heatmap de ejemplo de la estrategia en ejecución.}
    \label{fig:example_heatmap}
\end{figure}

\begin{figure}[htbp]
    \centering
    \includegraphics[width=0.8\linewidth]{./images/funding_rate/capital_growth.png}
    \caption{Evolución del capital invertido en la estrategia.}
    \label{fig:capital_growth}
\end{figure}

A continuación se muestran las ganancias anuales según datos históricos:

\begin{table}[htbp]
    \centering
    \begin{tabular}{|l|r|r|}
        \hline
        \textbf{Año} & \textbf{Ganancia (USDT)} & \textbf{\% sobre 2\,000 USDT} \\
        \hline
        2020 & 1\,136.38 & 56.82 \\
        2021 & 3\,063.87 & 153.19 \\
        2022 & 235.72  & 11.79 \\
        2023 & 710.63  & 35.53 \\
        2024 & 961.56  & 48.08 \\
        2025 & 145.37  & 7.27  \\
        \hline
    \end{tabular}
    \caption{Ganancias anuales de la estrategia de Funding Rate.}
    \label{tab:funding_rates_percent}
\end{table}

Por ejemplo, para el año 2024, el porcentaje de ganancia sobre el capital inicial es:
\[
\%\,\text{USDT} = \frac{961.56}{2\,000} \cdot 100 \approx 48.08\%
\]

\end{document}